\documentclass[11pt,a4paper]{article}
\usepackage[utf8]{inputenc}
\usepackage[margin=1in]{geometry}
\usepackage{amsmath}
\usepackage{amssymb}
\usepackage{enumitem}
\usepackage{booktabs}
\usepackage{titlesec}
\usepackage{microtype}
\usepackage{xcolor}

\titleformat{\section}{\large\bfseries\sffamily}{}{0em}{}[\titlerule]
\titleformat{\subsection}{\normalsize\bfseries\sffamily}{}{0em}{}

\setlength{\parindent}{0pt}
\setlength{\parskip}{8pt plus 2pt minus 1pt}

\title{\textbf{CAMS Exam: Comprehensive Missing Topics (50 Questions)}}
\author{\textbf{Based on Candidate Feedback from 2025-2026}}
\date{\today}

\begin{document}
	
	\maketitle
	
	\section*{IMPORTANT CANDIDATE FEEDBACK SUMMARY}
	\begin{itemize}
		\item \textbf{"FATF Immediate Outcomes (IOs) were heavily tested"}
		\item \textbf{"Questions on distinguishing FATF from FSRBs"}
		\item \textbf{"Specific USA PATRIOT sections — 314(a), 314(b), 319(a), 319(b)"}
		\item \textbf{"Lots of insurance (life insurance) scenarios"}
		\item \textbf{"Wolfsberg and Egmont secondary functions, not just primary"}
		\item \textbf{"Many questions required selecting 2-4 answers" }
		\item \textbf{"Extraterritorial reach of OFAC"}
	\end{itemize}
	
	\section*{Instructions}
	Each question has \textbf{five options (A through E)}. 
	The question will specify how many answers are correct: \textbf{[SELECT TWO]}, \textbf{[SELECT THREE]}, or \textbf{[SELECT FOUR]}.
	
	\newpage
	
	% ============================================================
	\section{FATF Immediate Outcomes (IO.1-IO.11)}
	% ============================================================
	
	\subsection{1. FATF Immediate Outcome 1 (IO.1) — Risk Understanding \textnormal{[SELECT TWO]}}
	
	A FATF assessor is evaluating a country's effectiveness under IO.1: "Money laundering and terrorist financing risks are understood and, where appropriate, actions coordinated domestically."
	
	Which two findings would most negatively affect the IO.1 rating?
	
	\begin{enumerate}[label=\Alph*.]
		\item The country completed a National Risk Assessment (NRA) five years ago but has not updated it despite significant changes in the financial sector and emergence of new virtual asset service providers.
		\item The financial supervisor has identified significant AML deficiencies at several banks but has not communicated these findings to the FIU or law enforcement.
		\item The country's FIU receives over 50,000 STRs annually and disseminates intelligence to law enforcement.
		\item The country has a formal inter-agency coordination committee that meets quarterly to discuss emerging ML/TF risks.
		\item The country's central bank publishes an annual AML/CFT supervision report on its website.
	\end{enumerate}
	
	\textbf{Correct Answers: A and B}
	
	\subsection{2. FATF Immediate Outcome 6 (IO.6) — Financial Intelligence Usage \textnormal{[SELECT TWO]}}
	
	IO.6 assesses whether "financial intelligence is used by competent authorities for ML/TF investigations."
	
	Which two findings would indicate a LOW effectiveness rating for IO.6?
	
	\begin{enumerate}[label=\Alph*.]
		\item The FIU receives 60,000 STRs annually but has never disseminated any intelligence to law enforcement.
		\item Law enforcement has opened 15 money laundering investigations in the past year, none of which were initiated by or used FIU intelligence.
		\item The FIU has a secure IT system for receiving and storing STRs.
		\item The FIU has a staff of 40 analysts who review all STRs within 30 days.
		\item Law enforcement and the FIU have a formal memorandum of understanding for sharing intelligence.
	\end{enumerate}
	
	\textbf{Correct Answers: A and B}
	
	\subsection{3. FATF Immediate Outcome 7 (IO.7) — ML Investigations \textnormal{[SELECT TWO]}}
	
	IO.7 assesses whether "money laundering is investigated and prosecuted."
	
	A country has excellent STR filing rates (80,000 annually) and a well-staffed FIU. However, law enforcement has opened only five ML investigations in four years, and there have been zero convictions. Prosecutors lack training on financial crime methodology.
	
	Which two statements correctly describe the likely IO.7 rating?
	
	\begin{enumerate}[label=\Alph*.]
		\item The IO.7 rating will likely be LOW because financial intelligence is not being used to initiate or support ML investigations, and the lack of convictions demonstrates systemic failure.
		\item The IO.7 rating will be HIGH because high STR volume alone demonstrates effectiveness regardless of investigation outcomes.
		\item The lack of prosecutor training on financial crime methodology is a material deficiency that will negatively affect the IO.7 rating.
		\item The existence of five investigations, even with zero convictions, is sufficient for a MODERATE rating.
		\item IO.7 only assesses the legal framework, not actual prosecution outcomes.
	\end{enumerate}
	
	\textbf{Correct Answers: A and C}
	
	\subsection{4. FATF Immediate Outcome 8 (IO.8) — Confiscation \textnormal{[SELECT TWO]}}
	
	IO.8 assesses whether "proceeds of crime and instrumentalities of crime are confiscated."
	
	A country has laws authorizing criminal and civil forfeiture. However, in five years, no assets have been confiscated because prosecutors lack training on tracing criminal proceeds, and there is no asset management office to handle seized property.
	
	Which two statements are correct?
	
	\begin{enumerate}[label=\Alph*.]
		\item The existence of confiscation laws alone does not satisfy IO.8; the country must demonstrate that proceeds are actually being identified, traced, frozen, and confiscated.
		\item The lack of prosecutor training on tracing and the absence of an asset management office are material deficiencies that will result in a LOW effectiveness rating for IO.8.
		\item As long as the confiscation laws exist on paper, IO.8 is automatically rated HIGH regardless of implementation.
		\item Confiscation only requires a conviction; asset tracing is not required for IO.8.
		\item Only criminal forfeiture counts for IO.8; civil forfeiture is not considered.
	\end{enumerate}
	
	\textbf{Correct Answers: A and B}
	
	\subsection{5. FATF Immediate Outcome 10 (IO.10) — TF Investigations \textnormal{[SELECT TWO]}}
	
	IO.10 assesses whether "terrorist financing is investigated and prosecuted."
	
	A country has never frozen any terrorist assets under UNSCR 1373 and has zero TF convictions in ten years, despite being a regional hub for a known terrorist group's fundraising activities.
	
	Which two statements correctly describe the deficiencies?
	
	\begin{enumerate}[label=\Alph*.]
		\item The failure to freeze terrorist assets under UNSCR 1373 is a direct violation of binding UN Security Council resolutions and will result in a LOW rating for IO.10.
		\item The zero TF convictions indicate that the country is not effectively identifying, investigating, or prosecuting TF activity, regardless of the legal framework.
		\item As long as the country has criminalized TF, IO.10 is automatically rated HIGH.
		\item TF investigations are optional; countries may choose to prioritize other crimes.
		\item UNSCR 1373 asset freezing obligations only apply to UN member states that have ratified the resolution; non-ratifying members are exempt.
	\end{enumerate}
	
	\textbf{Correct Answers: A and B}
	
	% ============================================================
	\section{FATF-Style Regional Bodies (FSRBs) vs. FATF}
	% ============================================================
	
	\subsection{6. Distinguishing FATF from FSRBs \textnormal{[SELECT TWO]}}
	
	A compliance officer is explaining the relationship between FATF and FATF-Style Regional Bodies (FSRBs) such as MONEYVAL, GAFILAT, APG, and MENAFATF.
	
	Which two statements correctly distinguish FATF from FSRBs?
	
	\begin{enumerate}[label=\Alph*.]
		\item FATF is a global inter-governmental body that sets the 40 Recommendations; FSRBs are autonomous regional bodies that promote, implement, and assess the execution of FATF Recommendations within their geographic regions.
		\item FSRBs conduct mutual evaluations of their member jurisdictions using the same methodology as FATF; FATF does not directly evaluate FSRB members unless they are also FATF members.
		\item FSRBs are subordinate operational branches staffed directly by FATF employees; they have no independent authority.
		\item FATF has the authority to impose sanctions on FSRBs that fail to comply with Recommendations; FSRBs have no reciprocal authority over FATF.
		\item FATF sets binding international law that all UN member states must follow; FSRBs are private consulting networks with no regulatory authority.
	\end{enumerate}
	
	\textbf{Correct Answers: A and B}
	
	\subsection{7. FSRB Mutual Evaluation Authority \textnormal{[SELECT TWO]}}
	
	A country is a member of a FATF-Style Regional Body (FSRB) but not a FATF member. The FSRB has scheduled a mutual evaluation of the country's AML/CFT framework.
	
	Which two statements correctly describe the FSRB's authority and the consequences of the evaluation?
	
	\begin{enumerate}[label=\Alph*.]
		\item FSRBs conduct mutual evaluations using the same FATF Methodology (Technical Compliance and Effectiveness across 11 Immediate Outcomes).
		\item The evaluation results are reported to FATF, and a country with significant deficiencies may be referred to FATF for potential public identification (Grey List or Black List).
		\item FSRBs have no authority to evaluate non-FATF members; only FATF can conduct mutual evaluations.
		\item FSRB evaluation results are confidential and cannot be shared with FATF.
		\item Countries that fail FSRB evaluations face automatic expulsion from the United Nations.
	\end{enumerate}
	
	\textbf{Correct Answers: A and B}
	
	% ============================================================
	\section{Egmont Group — Specific Functions}
	% ============================================================
	
	\subsection{8. Egmont Group: Beyond Definition \textnormal{[SELECT TWO]}}
	
	A compliance officer is reviewing the role of the Egmont Group in international AML cooperation.
	
	Which two statements correctly describe specific functions of the Egmont Group beyond its basic definition?
	
	\begin{enumerate}[label=\Alph*.]
		\item The Egmont Group provides a secure communication platform (Egmont Secure Web) for FIUs to exchange financial intelligence on a peer-to-peer basis.
		\item The Egmont Group has established the "Egmont Group Principles for Information Exchange Between FIUs," which govern how member FIUs must handle and protect shared intelligence.
		\item The Egmont Group has binding enforcement authority over member FIUs and can compel them to produce records.
		\item The Egmont Group can prosecute money laundering cases independently of national law enforcement.
		\item The Egmont Group sets the 40 Recommendations that countries must follow.
	\end{enumerate}
	
	\textbf{Correct Answers: A and B}
	
	\subsection{9. Egmont Group: Information Exchange Limitations \textnormal{[SELECT TWO]}}
	
	FIU Alpha has received intelligence from FIU Beta through the Egmont Group's secure platform. FIU Alpha wants to use this intelligence to support a criminal prosecution.
	
	Which two statements correctly describe the limitations on Egmont-exchanged intelligence?
	
	\begin{enumerate}[label=\Alph*.]
		\item Egmont-exchanged intelligence must be used only for AML/CFT purposes and cannot be used for tax investigations or commercial purposes without the originating FIU's consent.
		\item The receiving FIU must maintain confidentiality of the shared intelligence and cannot disclose it to third parties without the originating FIU's prior authorization.
		\item Egmont-exchanged intelligence is automatically admissible in criminal court without any further authentication.
		\item Egmont-exchanged intelligence can be sold to commercial data brokers for profit.
		\item The originating FIU has no control over how the receiving FIU uses the intelligence once it is shared.
	\end{enumerate}
	
	\textbf{Correct Answers: A and B}
	
	% ============================================================
	\section{USA PATRIOT Act Sections 319(a) and 319(b)}
	% ============================================================
	
	\subsection{10. USA PATRIOT Act Section 319(a) — Forfeiture of Funds \textnormal{[SELECT TWO]}}
	
	Federal law enforcement has identified that funds in a U.S. correspondent account are traceable to a foreign bank's involvement in money laundering.
	
	Which two statements correctly describe Section 319(a) of the USA PATRIOT Act?
	
	\begin{enumerate}[label=\Alph*.]
		\item Section 319(a) allows the U.S. government to seize funds from a correspondent account maintained for a foreign bank if the funds are traceable to money laundering or are the proceeds of a crime.
		\item Section 319(a) provides that funds seized from a correspondent account may be forfeited to the United States, subject to judicial process.
		\item Section 319(a) only applies to funds held in accounts for foreign central banks; commercial banks are exempt.
		\item Section 319(a) requires a criminal conviction before any funds can be seized from a correspondent account.
		\item Section 319(a) applies only to funds in accounts held by FATF Black List jurisdictions.
	\end{enumerate}
	
	\textbf{Correct Answers: A and B}
	
	\subsection{11. USA PATRIOT Act Section 319(b) — Record Production \textnormal{[SELECT TWO]}}
	
	A U.S. bank maintains a correspondent account for Foreign Bank X. Federal law enforcement serves a subpoena on the U.S. bank demanding records related to Foreign Bank X's account. Foreign Bank X refuses to comply.
	
	Which two statements correctly describe Section 319(b)?
	
	\begin{enumerate}[label=\Alph*.]
		\item Section 319(b) allows the U.S. government to demand that a foreign bank produce records related to its U.S. correspondent account.
		\item If the foreign bank refuses to comply with a subpoena, the U.S. bank may be required to terminate the correspondent account, and funds in the account may be seized.
		\item Section 319(b) only applies to banks located in FATF Black List jurisdictions.
		\item Foreign banks are never required to produce records to U.S. authorities; the U.S. bank is solely responsible for producing its own records.
		\item Section 319(b) requires the U.S. bank to pay the foreign bank's legal fees for producing records.
	\end{enumerate}
	
	\textbf{Correct Answers: A and B}
	
	\subsection{12. USA PATRIOT Act Section 319(a) vs. 319(b) \textnormal{[SELECT TWO]}}
	
	A compliance officer is explaining the difference between Sections 319(a) and 319(b) to a new analyst.
	
	Which two statements correctly distinguish these sections?
	
	\begin{enumerate}[label=\Alph*.]
		\item Section 319(a) addresses the forfeiture of funds in interbank accounts; Section 319(b) addresses the production of records from foreign banks.
		\item Section 319(a) allows seizure of funds traceable to money laundering; Section 319(b) requires foreign banks to maintain and produce records of their U.S. correspondent accounts.
		\item Both sections apply only to banks in FATF Black List jurisdictions.
		\item Section 319(a) requires a criminal conviction; Section 319(b) does not.
		\item Both sections exempt foreign central banks from all requirements.
	\end{enumerate}
	
	\textbf{Correct Answers: A and B}
	
	% ============================================================
	\section{Wolfsberg Group — Beyond CBDDQ}
	% ============================================================
	
	\subsection{13. Wolfsberg Group: Primary Purpose and Scope \textnormal{[SELECT TWO]}}
	
	A compliance officer is explaining the Wolfsberg Group's role to a new board member.
	
	Which two statements correctly describe the Wolfsberg Group's primary purpose and scope?
	
	\begin{enumerate}[label=\Alph*.]
		\item The Wolfsberg Group is a private association of 13 global banks that develops best practices and guidance for AML/CFT, particularly for correspondent banking and private banking.
		\item The Wolfsberg Group's guidance is not legally binding but is widely adopted as industry standard and referenced by regulators worldwide.
		\item The Wolfsberg Group is an inter-governmental body that sets legally binding AML standards for all UN member states.
		\item The Wolfsberg Group has the authority to impose sanctions on non-compliant banks.
		\item The Wolfsberg Group only provides guidance on sanctions screening, not on CDD or correspondent banking.
	\end{enumerate}
	
	\textbf{Correct Answers: A and B}
	
	\subsection{14. Wolfsberg Correspondent Banking Due Diligence Questionnaire (CBDDQ) \textnormal{[SELECT THREE]}}
	
	A correspondent bank is reviewing a completed CBDDQ from a prospective respondent bank.
	
	Which three areas are typically covered in the Wolfsberg CBDDQ?
	
	\begin{enumerate}[label=\Alph*.]
		\item The respondent bank's AML/CFT policies, procedures, and internal controls.
		\item The respondent bank's customer base, including high-risk segments like MSBs, NBFIs, and foreign PEPs.
		\item The respondent bank's independent testing schedule and results.
		\item The respondent bank's CEO's personal investment portfolio.
		\item The respondent bank's office furniture preferences.
	\end{enumerate}
	
	\textbf{Correct Answers: A, B, and C}
	
	% ============================================================
	\section{Basel Committee on Banking Supervision}
	% ============================================================
	
	\subsection{15. Basel Committee: AML/CFT Risk Management Integration \textnormal{[SELECT TWO]}}
	
	The Basel Committee on Banking Supervision has issued guidance on the "Sound Management of Risks Related to ML and TF."
	
	Which two statements correctly describe the Basel Committee's expectations?
	
	\begin{enumerate}[label=\Alph*.]
		\item AML/CFT risk management should be integrated into the bank's overall operational risk management framework across all three lines of defense.
		\item Banks should conduct Enterprise-Wide Risk Assessments (EWRAs) to identify ML/TF risks across all products, services, geographies, and customer types.
		\item AML compliance should be isolated as a standalone silo function reporting only to the Chief Compliance Officer, with no integration into broader risk management.
		\item The first line of defense is internal audit; the third line of defense is business operations.
		\item The Basel Committee's guidance applies only to banks in FATF Grey List jurisdictions.
	\end{enumerate}
	
	\textbf{Correct Answers: A and B}
	
	% ============================================================
	\section{Insurance AML (Life Insurance)}
	% ============================================================
	
	\subsection{16. Life Insurance CDD Requirements \textnormal{[SELECT TWO]}}
	
	A life insurance company is implementing its AML program. The compliance officer must determine when CDD applies.
	
	Which two statements correctly describe CDD triggers for life insurance products?
	
	\begin{enumerate}[label=\Alph*.]
		\item For life insurance policies where the premium exceeds USD/EUR 10,000 per year or the single premium exceeds USD/EUR 10,000, CDD must be applied.
		\item CDD must be applied to the policyholder and, once the beneficiary is designated, to the beneficiary as well.
		\item Life insurance policies are completely exempt from CDD requirements because they are not financial products.
		\item CDD is only required for life insurance policies with annual premiums exceeding USD/EUR 100,000.
		\item Beneficiaries of life insurance policies never require CDD, regardless of the policy value.
	\end{enumerate}
	
	\textbf{Correct Answers: A and B}
	
	\subsection{17. Life Insurance: Beneficiary Red Flags \textnormal{[SELECT TWO]}}
	
	A life insurance company is reviewing a policy with a \$5 million death benefit. The policyholder is a foreign PEP from a high-corruption jurisdiction. The named beneficiary is an offshore shell company with no verifiable business purpose. The premium is paid via a wire transfer from a bank in a secrecy haven.
	
	Which two statements correctly identify the AML red flags?
	
	\begin{enumerate}[label=\Alph*.]
		\item The combination of foreign PEP policyholder, offshore shell company beneficiary, and premium payment from a secrecy haven jurisdiction constitutes multiple red flags requiring Enhanced Due Diligence.
		\item The insurance company must identify the beneficial owner of the shell company beneficiary before policy issuance.
		\item Life insurance policies with foreign PEP policyholders are prohibited under all circumstances.
		\item The insurance company may ignore the shell company beneficiary because beneficiaries are not subject to CDD.
		\item Premiums from secrecy haven jurisdictions are automatically legitimate if the wire transfer is properly formatted.
	\end{enumerate}
	
	\textbf{Correct Answers: A and B}
	
	\subsection{18. Life Insurance: Surrender and Loan Features as ML Vehicles \textnormal{[SELECT TWO]}}
	
	A money launderer purchases a \$2 million life insurance policy with criminal proceeds, pays the premium in full, and then immediately takes a policy loan for 90\% of the cash surrender value. The loan proceeds are wired to an offshore account.
	
	Which two statements correctly describe how life insurance features can be exploited for money laundering?
	
	\begin{enumerate}[label=\Alph*.]
		\item The policy loan feature can be used as a layering technique, converting criminal proceeds deposited as premium into a "legitimate" policy loan check or wire transfer.
		\item The cash surrender value can be accessed through partial withdrawals or full surrender, extracting laundered funds from the insurance company.
		\item Life insurance policies cannot be used for money laundering because all insurance transactions are reported to FinCEN.
		\item Policy loans are always exempt from AML monitoring because they are not considered transactions.
		\item Insurance companies are not required to monitor policy loan activity under the BSA.
	\end{enumerate}
	
	\textbf{Correct Answers: A and B}
	
	% ============================================================
	\section{Broker-Dealer AML Obligations}
	% ============================================================
	
	\subsection{19. Broker-Dealer CDD Requirements \textnormal{[SELECT TWO]}}
	
	A broker-dealer is opening a new account for a corporate client. The compliance officer must determine the CDD requirements.
	
	Which two statements correctly describe broker-dealer CDD obligations?
	
	\begin{enumerate}[label=\Alph*.]
		\item Broker-dealers must identify and verify beneficial owners of legal entity customers at account opening.
		\item Broker-dealers must conduct ongoing monitoring of transactions for suspicious activity and file SARs as required.
		\item Broker-dealers are exempt from CDD requirements because they are regulated by the SEC, not FinCEN.
		\item Broker-dealers only need to verify the identity of the person opening the account; beneficial ownership identification is optional.
		\item Broker-dealers never need to file SARs because suspicious trading is reported to FINRA, not FinCEN.
	\end{enumerate}
	
	\textbf{Correct Answers: A and B}
	
	\subsection{20. Broker-Dealer: Red Flags in Securities Trading \textnormal{[SELECT TWO]}}
	
	A broker-dealer's AML monitoring system flags the following activity: a customer deposits \$500,000 in cash into a brokerage account, purchases 10,000 shares of a low-volume penny stock, and then immediately sells the shares at a 5\% loss to a related party account. The funds are then wired to an offshore bank account.
	
	Which two statements correctly identify the red flags?
	
	\begin{enumerate}[label=\Alph*.]
		\item The large cash deposit (structuring indicator) and immediate purchase and sale of a low-volume security at a loss to a related party is a classic market manipulation and layering pattern.
		\item The transaction may constitute "wash trading" or matched orders, where the customer is using the securities market to move funds between controlled accounts while creating the appearance of legitimate trading activity.
		\item The cash deposit is the only red flag; securities transactions are exempt from SAR filing.
		\item The 5\% loss confirms that the customer is a legitimate investor and the transaction is not suspicious.
		\item Broker-dealers are not required to monitor cash deposits because cash is not a security.
	\end{enumerate}
	
	\textbf{Correct Answers: A and B}
	
	% ============================================================
	\section{Money Services Business (MSB) Agents}
	% ============================================================
	
	\subsection{21. MSB Principal Responsibility for Agent Compliance \textnormal{[SELECT TWO]}}
	
	National Money Transfer Inc. is a licensed MSB with 5,000 agent locations. One agent location has processed \$2 million in wire transfers in a single month — 4,000\% above its historical average. All transfers are under \$900 and go to the same three countries. The agent's owner cannot explain the volume increase.
	
	Which two statements correctly describe the MSB's obligations?
	
	\begin{enumerate}[label=\Alph*.]
		\item The MSB is responsible for its agent network's compliance under BSA regulations and must investigate this anomaly.
		\item The MSB should consider filing a SAR on the agent location and potentially terminating the agent agreement.
		\item The MSB has no liability because each agent is an independent contractor responsible for its own compliance.
		\item The MSB should increase the agent's transaction limits to accommodate the new business volume.
		\item The MSB's only obligation is to provide the agent with a copy of its AML policy; no further action is required.
	\end{enumerate}
	
	\textbf{Correct Answers: A and B}
	
	\subsection{22. Unregistered MSB Operation \textnormal{[SELECT TWO]}}
	
	A convenience store is cashing payroll checks and personal checks for a fee, advertising "Check Cashing — No ID Required — Up to \$5,000." The store is not registered with FinCEN as an MSB.
	
	Which two statements correctly describe the violations?
	
	\begin{enumerate}[label=\Alph*.]
		\item The grocery store is required to register as an MSB with FinCEN because check cashing exceeding \$1,000 per person per day triggers MSB status.
		\item The store's failure to file CTRs for cash transactions over \$10,000 or SARs for suspicious patterns violates the Bank Secrecy Act.
		\item Grocery stores are explicitly exempt from MSB registration requirements under federal law regardless of check cashing volume.
		\item The store's "No ID Required" policy is fully compliant with BSA regulations.
		\item The store is only required to register if it cashes checks for non-customers; cashing checks for customers is exempt.
	\end{enumerate}
	
	\textbf{Correct Answers: A and B}
	
	% ============================================================
	\section{Prepaid Access Programs}
	% ============================================================
	
	\subsection{23. Prepaid Access: MSB Status and Reporting \textnormal{[SELECT TWO]}}
	
	A fintech company issues open-loop prepaid debit cards that can be loaded with cash at retail locations and used internationally. The cards can be reloaded up to \$15,000 annually per card.
	
	Which two statements correctly describe the AML obligations for the prepaid access program?
	
	\begin{enumerate}[label=\Alph*.]
		\item The issuer of prepaid access is considered a Money Services Business (MSB) and must register with FinCEN.
		\item The issuer must implement a risk-based AML program, including CDD on cardholders who exceed certain thresholds.
		\item Prepaid cards are completely exempt from all BSA requirements because they are not considered monetary instruments.
		\item The issuer has no obligation to monitor reload activity because each reload is a cash transaction below reporting thresholds.
		\item Prepaid cards cannot be used for cross-border currency movement because they are not physical currency.
	\end{enumerate}
	
	\textbf{Correct Answers: A and B}
	
	\subsection{24. Prepaid Cards: Structuring and Border Evasion \textnormal{[SELECT TWO]}}
	
	A traveler is stopped at customs carrying 50 prepaid debit cards, each loaded with \$9,500. The cards were purchased with cash at multiple retail locations over 30 days. The traveler is flying to a jurisdiction under FinCEN heightened scrutiny.
	
	Which two statements correctly describe the AML concerns?
	
	\begin{enumerate}[label=\Alph*.]
		\item The traveler may have violated structuring laws by purchasing each card at exactly \$9,500 to remain below the \$10,000 CTR threshold.
		\item The traveler may have evaded currency reporting requirements by using prepaid cards, which are classified as monetary instruments under BSA regulations, to transport the equivalent of \$475,000 in currency.
		\item Prepaid cards are retail products and are not subject to any currency controls, regardless of value or number.
		\item The only violation is a state consumer protection law violation for purchasing prepaid cards without a business license.
		\item The customs agency should simply confiscate the cards and release the traveler, as there is no AML reporting obligation.
	\end{enumerate}
	
	\textbf{Correct Answers: A and B}
	
	% ============================================================
	\section{Digital Assets: Wallet Types}
	% ============================================================
	
	\subsection{25. Custodial vs. Non-Custodial Wallets: AML Obligations \textnormal{[SELECT TWO]}}
	
	A regulated VASP offers both custodial wallets (where the VASP holds private keys) and non-custodial wallets (where the customer holds private keys).
	
	Which two statements correctly distinguish the AML obligations for each wallet type?
	
	\begin{enumerate}[label=\Alph*.]
		\item For custodial wallets, the VASP has full CDD obligations because it controls the funds and can identify the customer's transactions.
		\item For non-custodial wallets, the VASP's ability to conduct CDD is limited because the customer controls the private keys and the VASP cannot verify off-platform transactions.
		\item Non-custodial wallets are completely exempt from all AML obligations because the customer controls the private keys.
		\item Custodial wallets require no CDD because the VASP is only a technology provider, not a financial institution.
		\item Both wallet types have identical AML obligations; the distinction between custodial and non-custodial is irrelevant for AML purposes.
	\end{enumerate}
	
	\textbf{Correct Answers: A and B}
	
	\subsection{26. Privacy Coins: Technical Features and AML Risks \textnormal{[SELECT TWO]}}
	
	A compliance analyst is reviewing the AML risks associated with privacy coins such as Monero (XMR) and Zcash (ZEC).
	
	Which two technical features create money laundering risks for privacy coins?
	
	\begin{enumerate}[label=\Alph*.]
		\item Monero uses ring signatures and stealth addresses to obscure transaction trails, making it difficult to trace the origin and destination of funds.
		\item Zcash uses zero-knowledge proofs (zk-SNARKs) that allow selective disclosure of transaction details but can be configured for complete anonymity.
		\item Privacy coins are fully transparent and traceable on public blockchains, just like Bitcoin.
		\item Privacy coins are illegal in all jurisdictions globally and cannot be used by regulated VASPs.
		\item Privacy coins are exclusively issued and controlled by central banks, eliminating money laundering risk.
	\end{enumerate}
	
	\textbf{Correct Answers: A and B}
	
	% ============================================================
	\section{Cross-Border Data Transfer (GDPR vs. AML)}
	% ============================================================
	
	\subsection{27. GDPR vs. AML: Right to Erasure Exception \textnormal{[SELECT TWO]}}
	
	A European bank receives a GDPR "right to erasure" request from a customer demanding deletion of all personal data, including KYC documents and transaction records. The compliance officer notes that a SAR was filed on this customer two years ago.
	
	Which two statements correctly resolve the conflict?
	
	\begin{enumerate}[label=\Alph*.]
		\item GDPR Article 17(3) explicitly carves out exceptions to the right to erasure where retention is necessary for compliance with a legal obligation under EU or member state law.
		\item AML/CFT recordkeeping obligations constitute a legal obligation that overrides the right to erasure for records required to be retained (typically 5 years).
		\item The bank must delete all records immediately upon receiving a GDPR erasure request, regardless of AML retention requirements.
		\item The SAR filing must be deleted because GDPR applies to all personal data, including confidential law enforcement reports.
		\item The bank may delete the KYC documents but must retain the SAR filing; partial deletion satisfies both obligations.
	\end{enumerate}
	
	\textbf{Correct Answers: A and B}
	
	\subsection{28. Cross-Border Data Sharing: MLAT vs. GDPR \textnormal{[SELECT TWO]}}
	
	U.S. law enforcement has issued a subpoena to a European bank operating a branch in the U.S., demanding customer records for a money laundering investigation. The bank's European headquarters argues that GDPR prohibits transferring the data outside the EU.
	
	Which two statements correctly analyze the conflict?
	
	\begin{enumerate}[label=\Alph*.]
		\item The GDPR includes exceptions for cross-border data transfers that are necessary for important public interest grounds, including law enforcement cooperation.
		\item The bank must comply with the U.S. subpoena because the U.S. branch is subject to U.S. law; the bank should seek legal advice on GDPR compliance mechanisms such as Standard Contractual Clauses.
		\item GDPR absolutely prohibits any transfer of personal data to the U.S. under any circumstances, overriding all law enforcement requests.
		\item The bank should ignore the U.S. subpoena because European data protection laws take precedence over U.S. law enforcement requests.
		\item The bank must delete all customer data before responding to any U.S. subpoena to ensure GDPR compliance.
	\end{enumerate}
	
	\textbf{Correct Answers: A and B}
	
	% ============================================================
	\section{Corporate Transparency Act (CTA)}
	% ============================================================
	
	\subsection{29. Corporate Transparency Act: Reporting Requirements \textnormal{[SELECT TWO]}}
	
	A small business owner asks whether his LLC is required to report beneficial ownership information under the Corporate Transparency Act (CTA).
	
	Which two statements correctly describe CTA reporting obligations?
	
	\begin{enumerate}[label=\Alph*.]
		\item The reporting company (the LLC) is primarily responsible for filing beneficial ownership information with FinCEN.
		\item Reporting companies must file initial BOI reports within 90 days of formation (or 30 days for entities created after the effective date, depending on timing).
		\item The company's bank is responsible for filing BOI on behalf of the company; the company has no independent filing obligation.
		\item BOI reports are filed with the IRS, not FinCEN.
		\item Large operating companies with over 20 employees and \$5 million in gross receipts are exempt from reporting under the CTA.
	\end{enumerate}
	
	\textbf{Correct Answers: A and B}
	
	\subsection{30. CTA: Exempt Entities \textnormal{[SELECT TWO]}}
	
	A compliance officer is determining which types of entities are exempt from BOI reporting under the Corporate Transparency Act.
	
	Which two entities are EXEMPT from CTA reporting?
	
	\begin{enumerate}[label=\Alph*.]
		\item Large operating companies with more than 20 full-time employees in the U.S., a physical office in the U.S., and over \$5 million in gross receipts.
		\item Publicly traded companies that are subject to SEC reporting requirements.
		\item Every newly formed LLC regardless of size.
		\item Any company with a single owner is exempt from reporting.
		\item Non-profit organizations are required to report under the CTA with no exemptions.
	\end{enumerate}
	
	\textbf{Correct Answers: A and B}
	
	% ============================================================
	\section{FinCEN CDD Rule — Beneficial Ownership}
	% ============================================================
	
	\subsection{31. FinCEN CDD Rule: Two Prongs of UBO Identification \textnormal{[SELECT TWO]}}
	
	Under FinCEN's 2016 CDD Final Rule, covered financial institutions must identify beneficial owners of legal entity customers.
	
	Which two categories of individuals must be identified?
	
	\begin{enumerate}[label=\Alph*.]
		\item Any individual who directly or indirectly owns 25\% or more of the equity interests of the legal entity (Ownership Prong).
		\item At least one individual who exercises significant managerial control over the legal entity (Control Prong), such as a CEO, CFO, or other senior officer.
		\item All employees of the legal entity, regardless of their position.
		\item Any individual who has ever held a bank account in the entity's name.
		\item All shareholders regardless of their ownership percentage.
	\end{enumerate}
	
	\textbf{Correct Answers: A and B}
	
	\subsection{32. Beneficial Ownership: Aggregating Direct and Indirect Stakes \textnormal{[SELECT ONE]}}
	
	Ownership structure: Individual K owns 20\% of TargetCo directly. Individual K also owns 80\% of Holding LLC, which owns 15\% of TargetCo. The UBO threshold is 25\% combined.
	
	What is Individual K's total effective ownership and correct treatment?
	
	\begin{enumerate}[label=\Alph*.]
		\item 20\% — below threshold — no verification required.
		\item 32\% (20\% + (80\% × 15\%)) — exceeds threshold — must be identified and verified.
		\item 35\% (20\% + 15\% pass-through) — exceeds threshold — must be identified.
		\item Individual K does not qualify as a UBO because the indirect stake is held through a separate legal entity.
		\item The 25\% threshold applies only to direct ownership; indirect stakes are irrelevant.
	\end{enumerate}
	
	\textbf{Correct Answer: B}
	
	% ============================================================
	\section{PEP Definitions (Domestic vs. Foreign)}
	% ============================================================
	
	\subsection{33. PEP Classification: Domestic vs. Foreign \textnormal{[SELECT TWO]}}
	
	A private bank is onboarding two new clients: (1) a foreign head of state of a country with high corruption indices; (2) a domestic governor of the bank's home state.
	
	Which two statements correctly describe the PEP classification and required due diligence?
	
	\begin{enumerate}[label=\Alph*.]
		\item The foreign head of state is a Foreign PEP requiring Enhanced Due Diligence and senior management approval.
		\item The domestic governor is a Domestic PEP under FATF standards (if the country's laws define domestic PEPs), requiring EDD but potentially less intensive than foreign PEPs depending on national implementation.
		\item Neither client is a PEP because PEP status only applies to individuals with direct ownership of financial accounts.
		\item Both clients are International PEPs and require identical EDD procedures.
		\item Domestic PEPs are completely exempt from all enhanced due diligence under FATF standards.
	\end{enumerate}
	
	\textbf{Correct Answers: A and B}
	
	\subsection{34. PEP Family Members and Close Associates \textnormal{[SELECT TWO]}}
	
	A private bank is onboarding Ms. Z, who is the adult daughter of the Minister of Finance of a foreign country. Ms. Z holds no government position herself and wishes to deposit \$5 million from "family inheritance."
	
	Which two statements correctly describe her classification and obligations?
	
	\begin{enumerate}[label=\Alph*.]
		\item Ms. Z is classified as a PEP Family Member, requiring Enhanced Due Diligence and senior management approval.
		\item EDD for PEP family members requires establishing source of wealth and source of funds through verifiable, corroborative documentation.
		\item Ms. Z is not a PEP because she holds no government position, so standard CDD is sufficient.
		\item The ambiguous "family inheritance" from a minister's daughter is not a red flag.
		\item Only the Minister himself is a PEP; family members are never classified as PEPs.
	\end{enumerate}
	
	\textbf{Correct Answers: A and B}
	
	% ============================================================
	\section{Adverse Media Screening and OSINT}
	% ============================================================
	
	\subsection{35. Adverse Media: Sources and Actionable Findings \textnormal{[SELECT TWO]}}
	
	An EDD team discovers through OSINT that a corporate customer's CEO was named (not charged) in a foreign parliamentary corruption inquiry, and the customer's primary counterparty's director appears on an Interpol Red Notice.
	
	Which two statements correctly describe the required response?
	
	\begin{enumerate}[label=\Alph*.]
		\item The OSINT findings constitute credible adverse media requiring immediate escalation to the compliance committee.
		\item OSINT findings do not require a conviction to be actionable — reasonable suspicion of ML/TF risk, not proof of guilt, is the standard.
		\item Because no conviction has occurred, the adverse media is inconclusive and can be noted without further action.
		\item The Interpol Red Notice alone is sufficient to close the account immediately without further investigation.
		\item Adverse media screening only applies to sanctions lists; investigative journalism and government procurement records are not admissible as adverse media.
	\end{enumerate}
	
	\textbf{Correct Answers: A and B}
	
	\subsection{36. OSINT Sources: Credible vs. Non-Credible \textnormal{[SELECT TWO]}}
	
	A compliance analyst is conducting OSINT on a high-risk customer and has identified information from multiple sources.
	
	Which two sources are considered credible for AML adverse media screening?
	
	\begin{enumerate}[label=\Alph*.]
		\item Official government gazettes and sanctions lists.
		\item Reputable investigative journalism consortia (e.g., ICIJ, OCCRP).
		\item Unverified anonymous social media posts.
		\item Personal blog posts written by anonymous sources.
		\item Unmoderated online forums.
	\end{enumerate}
	
	\textbf{Correct Answers: A and B}
	
	% ============================================================
	\section{De-Risking Guidance (FATF and Wolfsberg)}
	% ============================================================
	
	\subsection{37. FATF Guidance on De-Risking \textnormal{[SELECT TWO]}}
	
	FATF has issued guidance addressing the de-risking phenomenon where financial institutions terminate entire categories of customer relationships indiscriminately.
	
	Which two statements are consistent with FATF's guidance?
	
	\begin{enumerate}[label=\Alph*.]
		\item Broad, categorical termination of entire classes of customer or geographic relationships is NOT the preferred risk management strategy under the risk-based approach.
		\item Financial institutions should conduct individual, documented risk assessments of respondent institutions rather than applying blanket exclusions based solely on geography or customer type.
		\item FATF requires banks to terminate all relationships with Grey List jurisdictions immediately.
		\item De-risking is explicitly required by FATF as the primary risk management tool.
		\item FATF prohibits all correspondent banking relationships with entities in any jurisdiction with an FATF rating below "Compliant."
	\end{enumerate}
	
	\textbf{Correct Answers: A and B}
	
	\subsection{38. Wolfsberg Group Guidance on De-Risking \textnormal{[SELECT TWO]}}
	
	The Wolfsberg Group has also issued guidance addressing de-risking in correspondent banking.
	
	Which two statements are consistent with Wolfsberg's guidance?
	
	\begin{enumerate}[label=\Alph*.]
		\item Banks should avoid indiscriminate de-risking that cuts off entire regions or categories of respondents without individual risk assessment.
		\item Where a relationship is terminated, the decision should be documented and based on a legitimate risk assessment, not solely on commercial considerations.
		\item Wolfsberg requires all banks to terminate relationships with any respondent located outside G20 countries.
		\item De-risking is the preferred method for managing correspondent banking risk under Wolfsberg principles.
		\item Wolfsberg prohibits any termination of correspondent banking relationships regardless of risk.
	\end{enumerate}
	
	\textbf{Correct Answers: A and B}
	
	% ============================================================
	\section{OFAC Extraterritorial Reach}
	% ============================================================
	
	\subsection{39. OFAC Primary vs. Secondary Sanctions \textnormal{[SELECT TWO]}}
	
	A European bank processes a EUR-denominated transaction for a non-U.S. client conducting business with an Iranian energy entity not individually listed on the OFAC SDN list. No U.S. dollars are used, no U.S. persons are involved.
	
	Which two statements correctly assess the bank's sanctions risk?
	
	\begin{enumerate}[label=\Alph*.]
		\item The bank faces U.S. secondary sanctions risk under the Iran Threat Reduction and Syria Human Rights Act (ITRSHRA) and related legislation, which can target non-U.S. entities conducting "significant" transactions with designated Iranian energy sector entities.
		\item The practical risk for the bank is that even if EU Blocking Statutes provide a defense in EU courts, U.S. regulators can impose secondary sanctions restricting the bank's USD correspondent banking access, effectively cutting it off from dollar-clearing markets.
		\item The bank has no U.S. sanctions risk because the transaction uses EUR and has no U.S. nexus; OFAC primary sanctions only apply to U.S. persons and U.S.-dollar transactions.
		\item The bank is automatically exempt under the EU Blocking Statute, which eliminates all U.S. secondary sanctions exposure.
		\item The bank's only risk is from EU sanctions, which do not restrict energy sector transactions with Iran.
	\end{enumerate}
	
	\textbf{Correct Answers: A and B}
	
	\subsection{40. OFAC 50\% Rule: Multi-Tier Calculation \textnormal{[SELECT ONE]}}
	
	Ownership chain: Individual X (SDN-listed) owns 55\% of Company A. Company A owns 40\% of Company B. Company B owns 60\% of Company C. Individual Y (non-sanctioned) owns 40\% of Company C directly.
	
	Is Company C a blocked entity under the OFAC 50\% Rule?
	
	\begin{enumerate}[label=\Alph*.]
		\item Yes, because the ownership chain traces back to an SDN-listed individual.
		\item No. Company A is blocked. Company B is NOT blocked because Company A's 40\% ownership is below 50\%. Because Company B is not blocked, its 60\% ownership of Company C does not transmit blocked status. Company C is not blocked.
		\item Company C is blocked because Company B's 60\% ownership exceeds 50\%, and Company B must be analyzed as blocked because it is owned by blocked Company A regardless of percentage.
		\item The 50\% rule does not apply to indirect ownership tiers; only direct ownership matters.
		\item Company C is automatically blocked because any entity in the chain of ownership of an SDN-listed person is blocked regardless of percentage.
	\end{enumerate}
	
	\textbf{Correct Answer: B}
	
	% ============================================================
	\section{High-Difficulty Mixed Scenarios}
	% ============================================================
	
	\subsection{41. FATF Recommendation 8: NPO Vulnerabilities \textnormal{[SELECT TWO]}}
	
	A non-profit organization operates in a conflict zone where terrorist groups are known to control humanitarian aid distribution. The NPO receives donations from foreign sources and makes large cash withdrawals.
	
	Which two statements correctly identify the NPO's vulnerabilities and the bank's obligations?
	
	\begin{enumerate}[label=\Alph*.]
		\item The NPO's operation in a conflict zone where terrorist groups control aid distribution creates a significant risk that funds may be diverted to terrorist financing.
		\item The bank must apply enhanced due diligence to the NPO account, including monitoring for cash withdrawals inconsistent with the stated charitable mission.
		\item NPOs operating in conflict zones are exempt from all AML/CFT requirements because of humanitarian access concerns.
		\item The bank should close the NPO account immediately without any investigation to avoid regulatory liability.
		\item FATF Recommendation 8 requires that all NPOs be banned from operating in conflict zones.
	\end{enumerate}
	
	\textbf{Correct Answers: A and B}
	
	\subsection{42. FATF Recommendation 16: Travel Rule for VASPs \textnormal{[SELECT ONE]}}
	
	A registered VASP receives a transfer of 3 Bitcoin (\$180,000) from another VASP with no originator or beneficiary information attached. The receiving VASP's compliance officer argues that because the customer has already passed KYC, the missing originator data is the sending VASP's problem.
	
	Why is this analysis incorrect?
	
	\begin{enumerate}[label=\Alph*.]
		\item The analysis is correct; the Travel Rule imposes obligations only on the originating VASP.
		\item The receiving VASP has an obligation to obtain the required originator and beneficiary information. If the information is not provided, the receiving VASP must take follow-up action, including requesting the information, delaying the transaction, and applying enhanced scrutiny or SAR filing if the information cannot be obtained.
		\item The Travel Rule applies only to fiat-to-crypto transactions; crypto-to-crypto transfers are exempt.
		\item The FATF Travel Rule threshold for virtual assets is \$1 million; this transfer is below threshold.
		\item The receiving VASP's only obligation is to verify the identity of the beneficiary; originator data is exclusively the sending VASP's responsibility.
	\end{enumerate}
	
	\textbf{Correct Answer: B}
	
	\subsection{43. FATF Recommendation 24: Bearer Shares \textnormal{[SELECT TWO]}}
	
	A bank is onboarding a corporate customer incorporated in a jurisdiction that still permits bearer shares. The customer provides a notarized statement that the bearer shares have been immobilized with a licensed custodian.
	
	Which two statements correctly describe the bank's obligations?
	
	\begin{enumerate}[label=\Alph*.]
		\item The bank must verify that the custodian can identify the bearer share holder at all times and that the jurisdiction prohibits future issuance of bearer shares.
		\item The bank still needs to identify the beneficial owner for CDD purposes; immobilization does not eliminate the need for UBO identification.
		\item The immobilization of bearer shares with a licensed custodian fully satisfies all beneficial ownership identification requirements.
		\item The bank should accept the customer's notarized statement without further inquiry into the bearer share holder's identity.
		\item The use of bearer shares presents no money laundering risk and requires no additional due diligence.
	\end{enumerate}
	
	\textbf{Correct Answers: A and B}
	
	\subsection{44. FATF Recommendation 25: Trusts \textnormal{[SELECT TWO]}}
	
	A discretionary trust with assets of \$15 million is opening an account. The trust has a corporate trustee, a protector (who can veto distributions), and beneficiaries identified only by class, not by name.
	
	Which two statements correctly describe the CDD obligations?
	
	\begin{enumerate}[label=\Alph*.]
		\item The bank must identify and verify the corporate trustee and its UBOs, the protector, and the settlor of the trust.
		\item The bank must identify beneficiaries to the extent they are determinable; for discretionary trusts where beneficiaries are not individually named, the bank must understand the class of beneficiaries and obtain information sufficient to identify them.
		\item Discretionary trusts are exempt from all CDD requirements because the trustee has full discretion over distributions.
		\item Only the corporate trustee requires identification; beneficiaries and protectors are never subject to CDD.
		\item Trusts established in Crown Dependencies are automatically eligible for Simplified Due Diligence.
	\end{enumerate}
	
	\textbf{Correct Answers: A and B}
	
	\subsection{45. FATF Recommendation 30: Financial Investigations \textnormal{[SELECT TWO]}}
	
	A country's law enforcement has specialized financial investigation units. However, these units cannot obtain financial records directly from banks; all record requests must go through the prosecutor's office, which takes 6-12 months.
	
	Which two statements correctly describe the deficiencies?
	
	\begin{enumerate}[label=\Alph*.]
		\item Law enforcement must have the power to obtain financial records expeditiously, without unreasonable delay.
		\item The 6-12 month delay in obtaining records undermines the effectiveness of financial investigations.
		\item A prosecutor-only approval system is acceptable as long as the law exists.
		\item Law enforcement does not need direct access to bank records; prosecutors are competent authorities.
		\item Financial investigations are optional; law enforcement can choose to investigate other crimes instead.
	\end{enumerate}
	
	\textbf{Correct Answers: A and B}
	
	\subsection{46. FATF Recommendation 32: Cash Courier Controls \textnormal{[SELECT TWO]}}
	
	A country has no declaration system for cross-border currency movement. Customs officers have intercepted 20 cash couriers carrying over \$15 million in undeclared currency in the past year, but no penalties are imposed because the law does not require declarations.
	
	Which two statements correctly describe the deficiencies?
	
	\begin{enumerate}[label=\Alph*.]
		\item Countries must have a declaration or disclosure system for cross-border transportation of currency and bearer negotiable instruments.
		\item Countries must have the authority to stop or restrain cash couriers and impose sanctions for false declarations or non-disclosure.
		\item Cash courier controls are optional; Recommendation 32 is only a best practice.
		\item Only cash amounts exceeding \$100,000 require declaration under FATF standards.
		\item Bearer negotiable instruments (traveler's checks, money orders) are exempt from declaration requirements.
	\end{enumerate}
	
	\textbf{Correct Answers: A and B}
	
	\subsection{47. FATF Recommendation 34: Guidance and Feedback \textnormal{[SELECT TWO]}}
	
	A country has not issued any AML/CFT guidance to financial institutions or provided any feedback on STR quality.
	
	Which two statements correctly describe the obligations under Recommendation 34?
	
	\begin{enumerate}[label=\Alph*.]
		\item Competent authorities should provide guidance and feedback to financial institutions to assist them in implementing AML/CFT measures effectively.
		\item Guidance may include red flag indicators, typologies, and emerging ML/TF trends.
		\item Guidance and feedback are optional; Recommendation 34 is a best practice with no compliance obligation.
		\item Feedback on STR quality can only be provided in person; written feedback is prohibited.
		\item Recommendation 34 only applies to non-profit organizations, not financial institutions.
	\end{enumerate}
	
	\textbf{Correct Answers: A and B}
	
	\subsection{48. FATF Immediate Outcome 5: Preventive Measures \textnormal{[SELECT TWO]}}
	
	IO.5 assesses whether "legal persons and arrangements are prevented from misuse for money laundering or terrorist financing."
	
	A country has no beneficial ownership registry and does not require companies to identify or update UBO information.
	
	Which two statements correctly describe the deficiencies?
	
	\begin{enumerate}[label=\Alph*.]
		\item The absence of a beneficial ownership registry makes it difficult for financial institutions to verify UBO information, undermining preventive measures.
		\item The lack of a requirement to update UBO information after initial collection means that information quickly becomes outdated and unreliable.
		\item IO.5 only applies to financial institutions, not to legal persons themselves.
		\item A beneficial ownership registry is optional under FATF standards; countries may choose not to maintain one.
		\item UBO information is only required for publicly traded companies; private companies are exempt.
	\end{enumerate}
	
	\textbf{Correct Answers: A and B}
	
	\subsection{49. FATF Immediate Outcome 2: International Cooperation \textnormal{[SELECT TWO]}}
	
	IO.2 assesses whether "international cooperation is used to obtain evidence from foreign jurisdictions."
	
	A country has never submitted or responded to an MLAT request and has not designated a central authority for MLAT processing.
	
	Which two statements correctly describe the deficiencies?
	
	\begin{enumerate}[label=\Alph*.]
		\item MLAT requests are the primary mechanism for obtaining bank records and evidence from foreign jurisdictions for criminal proceedings.
		\item The absence of a designated central authority for MLAT processing creates delays and inefficiencies in international cooperation.
		\item International cooperation is optional; countries may choose not to participate in MLATs.
		\item Egmont Group exchanges can replace MLATs for all purposes, including criminal prosecution.
		\item IO.2 only applies to extradition, not to evidence gathering.
	\end{enumerate}
	
	\textbf{Correct Answers: A and B}
	
	\subsection{50. Comprehensive Final: Multiple Entity Types \textnormal{[SELECT THREE]}}
	
	A compliance officer is reviewing the AML obligations for five different entity types operating in the same jurisdiction.
	
	Which three of these entities are considered financial institutions subject to AML/CFT obligations under FATF Recommendations?
	
	\begin{enumerate}[label=\Alph*.]
		\item Banks and credit unions.
		\item Money or value transfer services (MVTS) including MSBs.
		\item Casinos (with gross gaming revenue exceeding USD/EUR 1 million).
		\item Restaurants that accept cash payments.
		\item Convenience stores that sell prepaid cards but do not cash checks or transmit money.
	\end{enumerate}
	
	\textbf{Correct Answers: A, B, and C}
	
\end{document}